\section{Electromagnetism, Gravity, and the Quantum}

The gauge structure of the last section is a claim about symmetry. This one is about what the substrate actually measures when that symmetry is run, and it ranges over the three great classical pillars, the electromagnetic field, gravity, and the strange nonlocality of the quantum.

Electromagnetism leaves clean fingerprints, and the substrate shows them. Charge conservation, which the knit enforces line by line, becomes the Gauss law of the field, the statement that charge is the source of flux. In the flat cusp the static potential of a point charge falls off as $1/r$, the Coulomb law, measured directly. And the field is genuinely gauge, in the sense that the physical content lives in loops rather than in the potential itself, which the substrate confirms through a gauge-invariant Wilson loop and the Aharonov-Bohm phase a charge picks up encircling a flux (depth L2). These are not new physics, they are the known signatures of an emergent $U(1)$, and finding them is how we know the field is real rather than nominal.

\subsection{Gravity}

Gravity in this model is not a force carried by the tone but a property of the mesh itself, the way matter shapes the geometry it sits in. In the flat three-dimensional cusp this shows up first as a clean Newtonian potential, $1/r$, measured for a mass in the cusp (depth L2). Pushed to the effective field level the geometry obeys the Einstein equations, recovered along the lines Jacobson made familiar, as an equation of state relating the curvature to the flow of energy across local horizons, with gravitational radiation and the inspiral chirp following at the same level. The substrate, in other words, reproduces general relativity where general relativity has been tested, as an emergent description of how the cusp bends.

There is an honest gap, and it is named rather than hidden. Clean $1/r$ gravity is a result about the flat cusp of the committed substrate. It does not transfer cleanly to the curved bulk of a different honeycomb such as $\{5,3,4\}$, whose interior is hyperbolic and screens the potential and whose own cusp is only two-dimensional. Two attempts to extract a clean gravitational signal in that curved bulk did not discriminate at the sizes we can reach, because the bulk there is dominated by its boundary, so that case is reported as open rather than forced.

\subsection{Cosmology}

The cosmology needs no new mechanism, because the universe was already growing. The number of docks within a given depth of the bulk increases exponentially with that depth, and read as a scale factor in time this exponential is precisely a de Sitter expansion, $a(t) = e^{Ht}$, with a cosmological constant $\Lambda = 3H^2$ set by the growth rate. The accelerating expansion of space is the growing of the crystal seen from inside (depth L2). The same radial direction, depth into the bulk, doubles as the scale of a renormalization group, with each shell a coarser description of the one beneath it, so the holographic running of couplings and the cosmic expansion are one geometric fact viewed two ways.

\subsection{The quantum}

The hardest pillar for any local, deterministic substrate is the quantum, and specifically the violation of Bell's inequality. A theory whose base is strictly local, with each dock reading only its own sites, seems bound to respect the classical bound, a CHSH correlation of at most $2$. Yet nature, and the substrate, reach the quantum value of $2\sqrt{2}$, the Tsirelson bound. The resolution is geometric. The cusp is local, but the bulk beneath it is small across, joining distant cusp points by short paths through the depth, so what looks like spooky action at a distance on the three-dimensional surface is an ordinary short connection through the fourth dimension. The model needs, and has, an emergent holographic nonlocality to reach $2\sqrt{2}$, and it gets there without any nonlocal rule in the base (depth L3). The Born rule for probabilities and the reflection positivity that underlies a sensible quantum theory both follow from the reversibility of the beat, leaving the measurement problem itself, why one outcome rather than the superposition, as a genuine open question rather than a solved one.
